\section{第38届复赛第1题：平凸/平凹透镜曲面设计}
\subsection{题干}
宽束平行光正入射到折射率为 $n$ 的平凸透镜平面侧后汇聚于焦点 $F$；透镜顶点 $O$ 到 $F$ 的距离为 $r_0$，如图\ref{fig:cpho38q01}。求凸面形状。随后插入折射率为 $n'$、顶点到 $F$ 的距离为 $r_0'$ 的平凹透镜，使出射光重新平行，求凹面形状及入、出平行圆光束半径比。
\begin{figure}[htbp]\centering\begin{tikzpicture}[x=.9cm,y=.9cm,>=Stealth,font=\small]
\draw[dashed] (-3.5,0)--(3.2,0); \coordinate (F) at (3,0); \coordinate (O) at (0,0);
\draw[thick] (-.15,-1.5) .. controls (.75,-1.2) and (.75,1.2) .. (-.15,1.5); \draw[thick] (-.15,-1.5)--(-.15,1.5);
\foreach \y in {-1.15,0,1.15}{\draw[->] (-3.3,\y)--(-.15,\y); \draw[->] (-.1,\y)--(F);}
\fill (O) circle(1.4pt) node[below] {$O$}; \fill (F) circle(1.4pt) node[right] {$F$};
\draw[<->] (0,-.32)--(3,-.32) node[midway,below] {$r_0$}; \node[left] at(-.15,-1.3){$P$}; \node[above] at(.45,1.2){$n$};
\end{tikzpicture}
\caption{第38届复赛第1题平凸透镜会聚光路的TikZ重绘。}\label{fig:cpho38q01}\end{figure}
\subsection{参考解答整理}
\begin{figure}[htbp]\centering\begin{tikzpicture}[x=.62cm,y=.62cm,>=Stealth,font=\scriptsize]
\coordinate (F) at (0,0); \coordinate (O) at (3,0);
\draw[dashed] (-.5,0)--(5.1,0); \fill(F)circle(1.4pt)node[below]{$F$}; \fill(O)circle(1.4pt)node[below]{$O$};
\draw[thick] (3,-1.7)..controls(3.9,-1.2)and(3.9,1.2)..(3,1.7);
\draw[->] (3.05,1.25)--(F); \draw[->] (3.05,-1.25)--(F); \draw[->] (3.05,0)--(F);
\draw[<->] (0,-.45)--(3,-.45)node[midway,below]{$r_0$}; \node[above]at(3.55,1.15){$M$};
\end{tikzpicture}
\qquad
\begin{tikzpicture}[x=.62cm,y=.62cm,>=Stealth,font=\scriptsize]
\coordinate (F) at (0,0); \coordinate (Op) at (2.0,0); \coordinate (O) at (3.8,0);
\draw[dashed] (-.5,0)--(5.3,0); \fill(F)circle(1.4pt)node[below]{$F$}; \fill(Op)circle(1.4pt)node[below]{$O'$}; \fill(O)circle(1.4pt)node[below]{$O$};
\draw[thick] (2,-1.55)..controls(1.25,-1.1)and(1.25,1.1)..(2,1.55); \draw[thick] (2,-1.55)--(2,1.55);
\draw[->] (2.05,1.1)--(4.8,1.1); \draw[->] (2.05,-1.1)--(4.8,-1.1); \draw[->] (2.05,0)--(4.8,0);
\draw[<->] (0,-.45)--(2,-.45)node[midway,below]{$r_0'$}; \node[above]at(1.45,1.1){$M'$};
\end{tikzpicture}

\caption{第38届复赛第1题等光程推导的辅助图：左为平凸透镜曲面，右为平凹透镜凹面。}\label{fig:cpho38q01solution}\end{figure}
以 $F$ 为极点，有等光程关系
\[
r=n(r\cos\theta-r_0)+r_0,
\qquad r=\frac{(1-n)r_0}{1-n\cos\theta}.
\]
以从 $F$ 指向 $O$ 的方向为 $x$ 轴，凸面满足
\[
\boxed{\frac{\left(x-\frac{n}{n+1}r_0\right)^2}{\left(\frac{r_0}{n+1}\right)^2}-\frac{y^2}{\left(\frac{\sqrt{n^2-1}}{n+1}r_0\right)^2}=1.}
\]
因此凸面为旋转双曲面。对平凹透镜完全类似：
\[
r'=\frac{(1-n')r_0'}{1-n'\cos\theta},
\]
\[
\boxed{\frac{\left(x-\frac{n'}{n'+1}r_0'\right)^2}{\left(\frac{r_0'}{n'+1}\right)^2}-\frac{y^2}{\left(\frac{\sqrt{n'^2-1}}{n'+1}r_0'\right)^2}=1.}
\]
凹面也是旋转双曲面。若汇聚光最大夹角为 $\theta_{\max}$，半径比为
\[
\boxed{k=\frac{(1-n)(1-n'\cos\theta_{\max})r_0}{(1-n')(1-n\cos\theta_{\max})r_0'}.}
\]
\subsection{转写依据}
题干、图1a和公式依据本题题面转写稿、答案公式核对表和图件核对入口。图1a以TikZ重绘。本题第（1）问20分，第（2）问20分。
